\documentclass[11pt]{article}

% ============================================================
% PACKAGES
% ============================================================

\usepackage[margin=1in]{geometry}
\usepackage{amsmath,amssymb,amsfonts}
\usepackage{graphicx}
\usepackage{float}
\usepackage{subcaption}
\usepackage{booktabs}
\usepackage{hyperref}
\usepackage{xcolor}
\usepackage{enumitem}
\usepackage{microtype}

% ============================================================
% FORMATTING
% ============================================================

\setlength{\parindent}{0pt}
\setlength{\parskip}{8pt}

\hypersetup{
    colorlinks=true,
    linkcolor=blue,
    urlcolor=blue
}

\graphicspath{{./}}

% ============================================================
% DOCUMENT
% ============================================================

\begin{document}

\begin{center}
    {\LARGE \textbf{MA550: Homework 2}}\\[8pt]
    {\large Finite Differences and Immersed Boundary Methods}\\[10pt]
    Axel Fraud\\
    Colorado School of Mines\\
    October 2026
\end{center}

\vspace{0.5cm}

% ============================================================
% PROBLEM 1
% ============================================================

\section*{Problem 1: Finite-Difference Solution of the Poisson Equation}

\subsection*{Numerical Method}

We consider the two-dimensional Poisson equation

\begin{equation}
    \nabla^2 u(x,y)=f(x,y)
\end{equation}

on the square domain

\[
    \Omega=[0,2\pi]\times[0,2\pi].
\]

The Laplacian is approximated using the standard five-point
finite-difference stencil,

\begin{equation}
    \nabla^2 u(x_i,y_j)
    \approx
    \frac{
    u_{i+1,j}+u_{i-1,j}
    +u_{i,j+1}+u_{i,j-1}
    -4u_{i,j}
    }{\Delta x^2}.
\end{equation}

This approximation is second-order accurate in the grid spacing,

\[
    \nabla_h^2 u
    =
    \nabla^2u+\mathcal{O}(\Delta x^2).
\]

The resulting sparse linear system is written as

\begin{equation}
    A\mathbf{u}=\mathbf{f},
\end{equation}

where $A$ represents the discrete Laplacian.

\subsection*{Numerical Solution}

The numerical solution exhibits the expected oscillatory
structure of the prescribed solution.

Figure~\ref{fig:p1_solution} shows both the three-dimensional
surface and its two-dimensional representation.

\begin{figure}[H]
    \centering

    \begin{subfigure}{0.75\textwidth}
        \centering
        \includegraphics[width=\textwidth]{P1_3D.png}
        \caption{Three-dimensional numerical solution.}
    \end{subfigure}

    \vspace{0.3cm}

    \begin{subfigure}{0.65\textwidth}
        \centering
        \includegraphics[width=\textwidth]{P1_2D.png}
        \caption{Two-dimensional contour representation.}
    \end{subfigure}

    \caption{Finite-difference solution of the Poisson equation.}
    \label{fig:p1_solution}
\end{figure}

\subsection*{Convergence Analysis}

To evaluate numerical accuracy, the relative errors are
computed using the discrete $L_\infty$ and $L_2$ norms,

\begin{equation}
    E_\infty
    =
    \frac{
    \|u_h-u_{\mathrm{exact}}\|_\infty
    }{
    \|u_{\mathrm{exact}}\|_\infty
    },
\end{equation}

and

\begin{equation}
    E_2
    =
    \frac{
    \|u_h-u_{\mathrm{exact}}\|_2
    }{
    \|u_{\mathrm{exact}}\|_2
    }.
\end{equation}

The convergence results are shown in
Figure~\ref{fig:p1_convergence}.

\begin{figure}[H]
    \centering
    \includegraphics[width=0.9\textwidth]{P1_Convergence.png}
    \caption{
    Relative numerical errors as functions of grid spacing.
    The dashed reference line represents second-order
    convergence, $\mathcal{O}(\Delta x^2)$.
    }
    \label{fig:p1_convergence}
\end{figure}

Both error norms closely follow the second-order reference
line. This confirms that the implementation achieves the
expected convergence behavior of the five-point
finite-difference approximation.

% ============================================================
% PROBLEM 2
% ============================================================

\newpage

\section*{Problem 2: Poisson Equation with an Immersed Boundary}

\subsection*{Mathematical Formulation}

We extend the finite-difference Poisson solver by introducing
an immersed boundary inside the computational domain.

The governing equation becomes

\begin{equation}
    \nabla^2u=f+S^{ib}(q),
\end{equation}

where $S^{ib}$ is the immersed-boundary spreading operator
and $q$ represents the unknown boundary forcing.

The immersed-boundary constraint is imposed through the
interpolation operator,

\begin{equation}
    J^{ib}(u)=g.
\end{equation}

The immersed boundary is parameterized by

\begin{align}
    x^{ib}(\theta)&=\pi+\cos\theta,\\
    y^{ib}(\theta)&=\pi+\sin\theta,
\end{align}

where $\theta\in[0,2\pi]$.

\subsection*{Spreading and Interpolation}

The spreading operator transfers forces defined on the
immersed boundary to the Cartesian grid.

In continuous form,

\begin{equation}
    S^{ib}(q)(x,y)
    =
    \int_{\Gamma}
    q(s)\,
    \delta_a(x-X(s))
    \delta_a(y-Y(s))
    \,ds.
\end{equation}

The interpolation operator evaluates the Cartesian solution
at the immersed-boundary points,

\begin{equation}
    J^{ib}(u)(s)
    =
    \int_\Omega
    u(x,y)
    \delta_a(x-X(s))
    \delta_a(y-Y(s))
    \,dx\,dy.
\end{equation}

Here, $\delta_a$ denotes a regularized delta function.

After discretization, the coupled problem can be expressed as

\begin{equation}
    \begin{bmatrix}
        A & -S\\
        J & 0
    \end{bmatrix}
    \begin{bmatrix}
        \mathbf{u}\\
        \mathbf{q}
    \end{bmatrix}
    =
    \begin{bmatrix}
        \mathbf{f}\\
        \mathbf{g}
    \end{bmatrix}.
\end{equation}

\subsection*{Numerical Solution}

Figure~\ref{fig:p2_solution} illustrates the numerical
solution with the immersed boundary superimposed.

\begin{figure}[H]
    \centering

    \begin{subfigure}{0.70\textwidth}
        \centering
        \includegraphics[width=\textwidth]{P2_3D.png}
        \caption{Three-dimensional solution.}
    \end{subfigure}

    \vspace{0.3cm}

    \begin{subfigure}{0.55\textwidth}
        \centering
        \includegraphics[width=\textwidth]{P2_2D.png}
        \caption{Two-dimensional solution with immersed boundary.}
    \end{subfigure}

    \caption{
    Numerical solution of the Poisson equation
    with an immersed boundary.
    }
    \label{fig:p2_solution}
\end{figure}

\subsection*{Grid Convergence}

The convergence study is shown in
Figure~\ref{fig:p2_convergence}.

\begin{figure}[H]
    \centering
    \includegraphics[width=0.9\textwidth]{P2_Convergence.png}
    \caption{
    Relative errors for the immersed-boundary Poisson solver.
    The dashed line indicates second-order convergence.
    }
    \label{fig:p2_convergence}
\end{figure}

The numerical errors decrease as the Cartesian grid is
refined. The $L_\infty$ error closely follows the
second-order reference line, while the $L_2$ error also
decreases consistently.

\subsection*{Effect of Immersed-Boundary Spacing}

The spacing between immersed-boundary points influences
both numerical accuracy and computational cost.

Figures~\ref{fig:p2_spacing_accuracy} and
\ref{fig:p2_spacing_runtime} illustrate these effects.

\begin{figure}[H]
    \centering
    \includegraphics[width=0.7\textwidth]{P2_Spacing_Accuracy.png}
    \caption{
    Relative error as a function of immersed-boundary
    spacing normalized by Cartesian grid spacing.
    }
    \label{fig:p2_spacing_accuracy}
\end{figure}

The errors are largest when the immersed-boundary points
are spaced substantially closer than the Cartesian grid
spacing. Increasing the spacing initially produces a
significant improvement, after which the errors
approximately stabilize.

\begin{figure}[H]
    \centering
    \includegraphics[width=0.7\textwidth]{P2_Spacing_Runtime.png}
    \caption{
    Computational runtime as a function of
    immersed-boundary spacing.
    }
    \label{fig:p2_spacing_runtime}
\end{figure}

The computational cost decreases considerably as the
immersed-boundary spacing increases.

These results illustrate the trade-off between
immersed-boundary resolution and computational efficiency.
For this experiment, spacing comparable to the Cartesian
grid spacing provides a reasonable balance.

% ============================================================
% PROBLEM 3
% ============================================================

\section*{Problem 3: An Almost Minimal Surface}

\subsection*{Problem Formulation}

We now consider the homogeneous Poisson equation

\begin{equation}
    \nabla^2u=S^{ib}(q),
\end{equation}

with homogeneous Dirichlet conditions on the outer boundary,

\begin{equation}
    u(0,y)=u(2\pi,y)=u(x,0)=u(x,2\pi)=0.
\end{equation}

The immersed boundary is parameterized by

\begin{align}
    x^{ib}(\theta)&=\pi+\cos\theta,\\
    y^{ib}(\theta)&=\pi+\sin\theta,
\end{align}

and the prescribed values are

\begin{equation}
    u(x^{ib},y^{ib})
    =
    \cos(2x^{ib})\sin(2y^{ib}).
\end{equation}

Unlike the previous problem, the forcing term is zero
throughout the domain except for the contribution from
the immersed boundary.

The resulting solution is therefore determined by the
prescribed values along the circular boundary and the
homogeneous conditions on the outer boundary.

\subsection*{Numerical Results}

The computed solution is shown in
Figure~\ref{fig:p3_solution}.

\begin{figure}[H]
    \centering

    \begin{subfigure}{0.75\textwidth}
        \centering
        \includegraphics[width=\textwidth]{P3_3D.png}
        \caption{Three-dimensional almost minimal surface.}
    \end{subfigure}

    \vspace{0.3cm}

    \begin{subfigure}{0.65\textwidth}
        \centering
        \includegraphics[width=\textwidth]{P3_2D.png}
        \caption{Two-dimensional contour representation.}
    \end{subfigure}

    \caption{
    Numerical solution of the homogeneous Poisson equation
    with prescribed immersed-boundary values.
    }
    \label{fig:p3_solution}
\end{figure}

The surface exhibits positive and negative deformations
near the immersed boundary, corresponding to the
oscillatory prescribed boundary values.

Away from the immersed boundary, the solution gradually
approaches zero as required by the homogeneous
outer Dirichlet conditions.

The resulting surface resembles a minimal surface in
the sense that it smoothly connects the prescribed
boundary values. However, it is not a true minimal
surface because the governing equation is linear,
whereas the minimal-surface equation is nonlinear.

% ============================================================
% PROBLEM 4: CHARLIE'S RULE
% ============================================================

\section*{Problem 4: Charlie's Rule}

For this problem, I explored several extensions of the immersed
boundary method developed in Problem 3. In particular, I investigated
how multiple moving immersed boundaries influence the numerical
solution and whether reusing previous solutions as initial guesses
could improve the efficiency of GMRES.

The goal was to extend the original static immersed boundary
problem into a pseudo-dynamic simulation while examining both
the resulting surface geometries and computational performance.

% ------------------------------------------------------------
% MOVING IMMERSED BOUNDARIES
% ------------------------------------------------------------

\subsection*{Moving Immersed Boundaries}

Rather than considering a single stationary immersed boundary,
I introduced two immersed boundaries whose positions and geometries
vary throughout the simulation.

The first boundary is an ellipse that undergoes rotation,
while the second is a circle that undergoes periodic deformation.

At each frame, the immersed boundary locations are updated,
and the corresponding numerical solution is recomputed.

This produces a pseudo-dynamic simulation in which the geometry
changes over time, but the governing elliptic problem is solved
independently at each configuration.

Mathematically, the numerical solution at frame $k$ satisfies

\begin{equation}
    \mathcal{L}_h u^{(k)}
    +
    \mathcal{S}^{(k)} F^{(k)}
    =
    f,
\end{equation}

where $\mathcal{L}_h$ is the discrete Laplacian,
$\mathcal{S}^{(k)}$ is the spreading operator associated
with the immersed boundary geometry at frame $k$, and
$F^{(k)}$ represents the immersed boundary forces.

The dependence on $k$ reflects the changing geometry
of the immersed boundaries.

% ------------------------------------------------------------
% 2D VISUALIZATION
% ------------------------------------------------------------

\subsection*{Two-Dimensional Visualization}

Figure~\ref{fig:p4_2d} shows two representative frames
from the two-dimensional simulation.

The black curves represent the immersed boundaries,
while the background contour plot illustrates the
numerical solution $u(x,y)$.

As the boundaries move and deform, the surrounding
solution changes accordingly.

In particular, the interaction between the two
immersed boundaries produces increasingly complex
spatial structures that would not appear in the
single-boundary configuration.

\begin{figure}[H]
    \centering

    \begin{subfigure}[b]{0.48\textwidth}
        \centering
        \includegraphics[width=\linewidth]{P4-2D-F1.png}
        \caption{Initial configuration.}
        \label{fig:p4_2d_f1}
    \end{subfigure}
    \hfill
    \begin{subfigure}[b]{0.48\textwidth}
        \centering
        \includegraphics[width=\linewidth]{P4-2D-F2.png}
        \caption{Later configuration.}
        \label{fig:p4_2d_f2}
    \end{subfigure}

    \caption{
    Two-dimensional snapshots of the moving immersed
    boundary simulation. The ellipse rotates while
    the circular boundary deforms, producing changes
    in the numerical solution.
    }
    \label{fig:p4_2d}
\end{figure}

% ------------------------------------------------------------
% 3D VISUALIZATION
% ------------------------------------------------------------

\subsection*{Three-Dimensional Surface Deformation}

To better visualize the influence of the immersed
boundaries, I also generated three-dimensional
representations of the numerical solution.

Figure~\ref{fig:p4_3d} shows two representative
frames from the resulting animation.

The surface height corresponds to the numerical
solution $u(x,y)$, while the black curves indicate
the immersed boundaries embedded within the surface.

As the boundaries evolve, the resulting surface
develops changing peaks, valleys, and localized
deformations.

These visualizations demonstrate how relatively
simple geometric motions can generate complicated
surface structures through the immersed boundary method.

\begin{figure}[H]
    \centering

    \begin{subfigure}[b]{0.48\textwidth}
        \centering
        \includegraphics[width=\linewidth]{P4-3D-F1.png}
        \caption{First surface configuration.}
        \label{fig:p4_3d_f1}
    \end{subfigure}
    \hfill
    \begin{subfigure}[b]{0.48\textwidth}
        \centering
        \includegraphics[width=\linewidth]{P4-3D-F2.png}
        \caption{Later surface configuration.}
        \label{fig:p4_3d_f2}
    \end{subfigure}

    \caption{
    Three-dimensional snapshots illustrating the
    deformation of the numerical solution as the
    immersed boundaries move. The changing boundary
    configurations produce different localized
    surface features.
    }
    \label{fig:p4_3d}
\end{figure}

% ------------------------------------------------------------
% GMRES PERFORMANCE
% ------------------------------------------------------------

\subsection*{GMRES Performance: Warm versus Cold Starts}

An additional experiment investigated whether
reusing the solution from the previous frame
could improve the convergence of GMRES.

In the original implementation, GMRES begins
with a zero initial guess,

\begin{equation}
    u_0^{(k)} = 0.
\end{equation}

This is referred to as a \textit{cold start}.

However, because the immersed boundaries move
continuously between frames, the numerical
solutions at consecutive frames are expected
to be similar.

This suggests using the previous solution
as the initial guess for the next solve,

\begin{equation}
    u_0^{(k)} = u^{(k-1)},
\end{equation}

which is referred to as a \textit{warm start}.

The motivation is that if

\begin{equation}
    \left\|
    u^{(k)} - u^{(k-1)}
    \right\|
    \ll
    \left\|u^{(k)}\right\|,
\end{equation}

then the previous solution should provide
a better approximation to the new solution
than the zero vector.

In principle, this could reduce the number
of Krylov iterations required for convergence.

\begin{figure}[H]
    \centering
    \includegraphics[width=0.90\textwidth]
    {P4_GMRES_Performance.png}

    \caption{
    Comparison of GMRES performance using cold
    and warm starts over successive animation
    frames. The upper panel shows the number
    of GMRES iterations required for convergence,
    while the lower panel compares the corresponding
    solution times.
    }
    \label{fig:p4_gmres}
\end{figure}

\subsection*{Discussion of GMRES Results}

Interestingly, the numerical results indicate
that warm-starting GMRES does not provide a
significant computational advantage for this
particular experiment.

As shown in Figure~\ref{fig:p4_gmres},
both methods require approximately 50--54
iterations across the simulation.

Furthermore, the solution times remain
comparable, with the warm-start approach
occasionally requiring slightly more time
than the cold-start approach.

Although the previous solution provides
a physically reasonable initial guess,
this does not necessarily guarantee faster
convergence of the Krylov subspace method.

GMRES convergence depends on several factors,
including the structure of the linear system,
the initial residual, and the spectral
properties of the operator.

Consequently, a better initial approximation
does not always produce a substantial reduction
in the number of iterations.

These results suggest that for this particular
system, improving the preconditioning strategy
may be more effective than simply reusing
the previous solution.

% ------------------------------------------------------------
% CONCLUSION
% ------------------------------------------------------------

\subsection*{Conclusion}

This experiment demonstrates how the immersed
boundary method can be extended beyond a
single stationary boundary to simulate
multiple moving and deforming geometries.

The resulting two-dimensional and
three-dimensional visualizations illustrate
how changes in boundary geometry influence
the numerical solution.

Additionally, the GMRES comparison provides
an interesting computational observation:
although warm-starting is a natural strategy
for a sequence of related linear systems,
it does not necessarily improve convergence
or computational efficiency.

Overall, these experiments highlight both
the flexibility of the immersed boundary
method and the importance of testing
numerical optimization strategies rather
than assuming they will improve performance.

\end{document}